\documentclass[10pt]{article}
% Shared preamble of the problems/MHD/setup_notes documents.
% Each case document does % Shared preamble of the problems/MHD/setup_notes documents.
% Each case document does % Shared preamble of the problems/MHD/setup_notes documents.
% Each case document does \input{preamble} right after \documentclass.
\usepackage[margin=0.85in]{geometry}
\usepackage[T1]{fontenc}
\usepackage{lmodern}
\usepackage{amsmath,amssymb,bm}
\usepackage{array,booktabs,tabularx,longtable,multirow}
\usepackage{xcolor}
\usepackage{enumitem}
\usepackage[normalem]{ulem}
\usepackage[hyphens]{url}
\usepackage[colorlinks=true,linkcolor=blue!50!black,urlcolor=blue!50!black]{hyperref}
\usepackage{caption}
\usepackage{graphicx}
\setlist{itemsep=1pt,topsep=3pt}
\sloppy
\emergencystretch=3em
\newcolumntype{Y}{>{\raggedright\arraybackslash}X}
\newcolumntype{L}[1]{>{\raggedright\arraybackslash}p{#1}}

% \code{...}: typewriter with breakable underscores (write \_ in the argument).
\makeatletter
\newcommand{\code}[1]{{\def\_{\char`\_\discretionary{}{}{}}\texttt{#1}}}
\makeatother

% --- Common-item marks -------------------------------------------------------
% \cm{x}  solid underline : identical in EVERY case to which the item applies
% \mj{x}  dashed underline: shared by a MAJORITY of those cases, but not all
\newcommand{\cm}[1]{\uline{#1}}
\newcommand{\mj}[1]{\dashuline{#1}}

\newcommand{\commonlegend}{%
\begin{center}
\fcolorbox{black!40}{black!3}{\parbox{0.92\textwidth}{\small
\textbf{How to read the underlines.}
\cm{Solid underline}: the setting (or the code path) is \emph{identical in every case
of \code{problems/MHD} to which it applies} --- all nine cases for general settings, the eight
\code{DSGS\_MHD} cases for DynSGS settings, the three \code{:positivity\_method => "conservative"}
cases for the element limiter.
\mj{Dashed underline}: shared by a \emph{majority} of those cases but not all of them.
No underline: specific to this case or to a minority.
The cross-case comparison is in \code{summary.pdf} of this directory.}}
\end{center}}

% --- Notation ----------------------------------------------------------------
\newcommand{\bq}{\mathbf{q}}
\newcommand{\bv}{\mathbf{v}}
\newcommand{\bmom}{\mathbf{m}}
\newcommand{\bB}{\mathbf{B}}
\newcommand{\bF}{\mathbf{F}}
\newcommand{\Et}{\tilde E}
\newcommand{\eps}{\varepsilon}
\newcommand{\dsgs}{DynSGS}
\newcommand{\CR}{C_R}
\newcommand{\Cmax}{C_{\max}}
\newcommand{\Cmin}{C_{\min}}
\newcommand{\nuc}{\nu_{c}}
\newcommand{\nufl}{\nu_{\mathrm{fl}}}
\newcommand{\commit}{\code{dcc3653}}
\newcommand{\branch}{\code{sm/MHD}}

% Two-column "setting" table: key | value as used | where / remark
\newenvironment{keytable}{%
\par\smallskip\noindent\small
\begin{longtable}{@{}L{0.24\textwidth}L{0.38\textwidth}L{0.32\textwidth}@{}}
\toprule
\textbf{Deck key / item} & \textbf{Value as run} & \textbf{Source / remark}\\
\midrule\endhead
\bottomrule\endfoot}{%
\end{longtable}\par}
 right after \documentclass.
\usepackage[margin=0.85in]{geometry}
\usepackage[T1]{fontenc}
\usepackage{lmodern}
\usepackage{amsmath,amssymb,bm}
\usepackage{array,booktabs,tabularx,longtable,multirow}
\usepackage{xcolor}
\usepackage{enumitem}
\usepackage[normalem]{ulem}
\usepackage[hyphens]{url}
\usepackage[colorlinks=true,linkcolor=blue!50!black,urlcolor=blue!50!black]{hyperref}
\usepackage{caption}
\usepackage{graphicx}
\setlist{itemsep=1pt,topsep=3pt}
\sloppy
\emergencystretch=3em
\newcolumntype{Y}{>{\raggedright\arraybackslash}X}
\newcolumntype{L}[1]{>{\raggedright\arraybackslash}p{#1}}

% \code{...}: typewriter with breakable underscores (write \_ in the argument).
\makeatletter
\newcommand{\code}[1]{{\def\_{\char`\_\discretionary{}{}{}}\texttt{#1}}}
\makeatother

% --- Common-item marks -------------------------------------------------------
% \cm{x}  solid underline : identical in EVERY case to which the item applies
% \mj{x}  dashed underline: shared by a MAJORITY of those cases, but not all
\newcommand{\cm}[1]{\uline{#1}}
\newcommand{\mj}[1]{\dashuline{#1}}

\newcommand{\commonlegend}{%
\begin{center}
\fcolorbox{black!40}{black!3}{\parbox{0.92\textwidth}{\small
\textbf{How to read the underlines.}
\cm{Solid underline}: the setting (or the code path) is \emph{identical in every case
of \code{problems/MHD} to which it applies} --- all nine cases for general settings, the eight
\code{DSGS\_MHD} cases for DynSGS settings, the three \code{:positivity\_method => "conservative"}
cases for the element limiter.
\mj{Dashed underline}: shared by a \emph{majority} of those cases but not all of them.
No underline: specific to this case or to a minority.
The cross-case comparison is in \code{summary.pdf} of this directory.}}
\end{center}}

% --- Notation ----------------------------------------------------------------
\newcommand{\bq}{\mathbf{q}}
\newcommand{\bv}{\mathbf{v}}
\newcommand{\bmom}{\mathbf{m}}
\newcommand{\bB}{\mathbf{B}}
\newcommand{\bF}{\mathbf{F}}
\newcommand{\Et}{\tilde E}
\newcommand{\eps}{\varepsilon}
\newcommand{\dsgs}{DynSGS}
\newcommand{\CR}{C_R}
\newcommand{\Cmax}{C_{\max}}
\newcommand{\Cmin}{C_{\min}}
\newcommand{\nuc}{\nu_{c}}
\newcommand{\nufl}{\nu_{\mathrm{fl}}}
\newcommand{\commit}{\code{dcc3653}}
\newcommand{\branch}{\code{sm/MHD}}

% Two-column "setting" table: key | value as used | where / remark
\newenvironment{keytable}{%
\par\smallskip\noindent\small
\begin{longtable}{@{}L{0.24\textwidth}L{0.38\textwidth}L{0.32\textwidth}@{}}
\toprule
\textbf{Deck key / item} & \textbf{Value as run} & \textbf{Source / remark}\\
\midrule\endhead
\bottomrule\endfoot}{%
\end{longtable}\par}
 right after \documentclass.
\usepackage[margin=0.85in]{geometry}
\usepackage[T1]{fontenc}
\usepackage{lmodern}
\usepackage{amsmath,amssymb,bm}
\usepackage{array,booktabs,tabularx,longtable,multirow}
\usepackage{xcolor}
\usepackage{enumitem}
\usepackage[normalem]{ulem}
\usepackage[hyphens]{url}
\usepackage[colorlinks=true,linkcolor=blue!50!black,urlcolor=blue!50!black]{hyperref}
\usepackage{caption}
\usepackage{graphicx}
\setlist{itemsep=1pt,topsep=3pt}
\sloppy
\emergencystretch=3em
\newcolumntype{Y}{>{\raggedright\arraybackslash}X}
\newcolumntype{L}[1]{>{\raggedright\arraybackslash}p{#1}}

% \code{...}: typewriter with breakable underscores (write \_ in the argument).
\makeatletter
\newcommand{\code}[1]{{\def\_{\char`\_\discretionary{}{}{}}\texttt{#1}}}
\makeatother

% --- Common-item marks -------------------------------------------------------
% \cm{x}  solid underline : identical in EVERY case to which the item applies
% \mj{x}  dashed underline: shared by a MAJORITY of those cases, but not all
\newcommand{\cm}[1]{\uline{#1}}
\newcommand{\mj}[1]{\dashuline{#1}}

\newcommand{\commonlegend}{%
\begin{center}
\fcolorbox{black!40}{black!3}{\parbox{0.92\textwidth}{\small
\textbf{How to read the underlines.}
\cm{Solid underline}: the setting (or the code path) is \emph{identical in every case
of \code{problems/MHD} to which it applies} --- all nine cases for general settings, the eight
\code{DSGS\_MHD} cases for DynSGS settings, the three \code{:positivity\_method => "conservative"}
cases for the element limiter.
\mj{Dashed underline}: shared by a \emph{majority} of those cases but not all of them.
No underline: specific to this case or to a minority.
The cross-case comparison is in \code{summary.pdf} of this directory.}}
\end{center}}

% --- Notation ----------------------------------------------------------------
\newcommand{\bq}{\mathbf{q}}
\newcommand{\bv}{\mathbf{v}}
\newcommand{\bmom}{\mathbf{m}}
\newcommand{\bB}{\mathbf{B}}
\newcommand{\bF}{\mathbf{F}}
\newcommand{\Et}{\tilde E}
\newcommand{\eps}{\varepsilon}
\newcommand{\dsgs}{DynSGS}
\newcommand{\CR}{C_R}
\newcommand{\Cmax}{C_{\max}}
\newcommand{\Cmin}{C_{\min}}
\newcommand{\nuc}{\nu_{c}}
\newcommand{\nufl}{\nu_{\mathrm{fl}}}
\newcommand{\commit}{\code{dcc3653}}
\newcommand{\branch}{\code{sm/MHD}}

% Two-column "setting" table: key | value as used | where / remark
\newenvironment{keytable}{%
\par\smallskip\noindent\small
\begin{longtable}{@{}L{0.24\textwidth}L{0.38\textwidth}L{0.32\textwidth}@{}}
\toprule
\textbf{Deck key / item} & \textbf{Value as run} & \textbf{Source / remark}\\
\midrule\endhead
\bottomrule\endfoot}{%
\end{longtable}\par}

\title{\code{problems/MHD/brioWu1d}\\[2pt]
\large DynSGS and positivity setup, as coded on branch \branch{} (commit \commit)}
\author{}\date{}
\begin{document}
\maketitle
\commonlegend

\noindent\textbf{Scope.} Everything below is read off the deck (\code{user\_inputs.jl}), the
case files and the kernels they reach on \branch{}; defaults are the ones \code{mod\_inputs.jl}
fills in when the deck is silent, and environment overrides are taken at their unset values
(\code{JEXPRESSO\_BW\_*}). Comments in the files were not taken as evidence where the code says
otherwise; the disagreements found are listed in Section~\ref{sec:obs}.

\section{The problem as set up}
\begin{keytable}
Equations & 1D ideal MHD, \cm{state $(\rho,\rho u,\rho v,E,\rho w,B_x,B_y,B_z)$}, 8 slots, no GLM field & \code{initialize.jl}, \code{user\_flux.jl}\\
EOS & \cm{$p=(\gamma-1)(E-\tfrac12|\bmom|^2/\rho-\tfrac12|\bB|^2)$}, $\gamma = 2$ & \code{user\_flux.jl} ($\gamma_{mhd}$\code{ = 2.0})\\
Domain, mesh & $x\in(0,1)$, 150 elements (600 DOFs) & \code{:nelx} $=600/N$ (\code{JEXPRESSO\_BW\_DOFS}, \code{\_NELX})\\
Discretization & \cm{CG spectral elements, LGL nodes, lumped (diagonal) mass} & \code{:interpolation\_nodes => "lgl"}\\
Polynomial order & \mj{$N = 4$} & \code{:nop} (\code{JEXPRESSO\_BW\_NOP})\\
Time integration & \cm{\code{CarpenterKennedy2N54}}, fixed \cm{$\Delta t$}, $\Delta t = 5\times10^{-5}$, $t_{end}=0.1$ & $\Delta t = 5\times10^{-5}\cdot600/(n_{el}N)$\\
Variables & \cm{\code{TOTAL()}}, \cm{\code{:energy\_equation => "energy"}} & \\
Initial data & left $(\rho,u,p,B_y)=(1,0,1,1)$, right $(0.125,0,0.1,-1)$, $B_x=0.75$, jump at $x=0.5$ & \code{initialize.jl}\\
Boundary conditions & Dirichlet on all 8 slots at both ends (the left/right states) & \code{user\_bc.jl}\\
Sources & none (\code{:lsource => false}) & \code{user\_source.jl}\\
Other switches & \cm{\code{:lfilter}, \code{:lkep}, \code{:entropy\_variables}, \code{:lrichardson}, \code{:ladapt} all false} & deck / defaults\\
\end{keytable}

\section{DynSGS: deck keys and effective values}
The deck selects \cm{\code{:visc\_model => DSGS\_MHD()}} with \cm{\code{:lvisc => true}}, which in 1D
reaches \code{compute\_dsgs\_viscosity!(::DSGS\_MHD, ::NSD\_1D)} (\code{SGS.jl:1929}) through
\code{viscous\_rhs\_el!(\dots, ::NSD\_1D)} (\code{rhs.jl:1398}).

The deck rescales the three constants by $h_{fac}=(N+1)/N = 5/4$
(\code{JEXPRESSO\_BW\_HSCALE} unset $\Rightarrow$ \code{"nop"}): the kernel uses
$\Delta = \Delta x/(N+1)$, and the scaling turns it into Dao \& Nazarov's $\Delta x/N$.
\begin{keytable}
\code{:visc\_model} & \cm{\code{DSGS\_MHD()}} & \\
\code{:}$\mu$ (slot multipliers) & \cm{all 1} (8 entries) & mass diffusion is on\\
\code{:dsgs\_sensor} & \code{"residual"} & element residual, stage-consistent stencil\\
\code{:dsgs\_CR} & $1.0\cdot h_{fac}^2 = 1.5625$ & $\equiv C_R=1$ with $\Delta x/N$\\
\code{:dsgs\_Cmax} & $0.5\cdot h_{fac} = 0.625$ & $\equiv C_{\max}=0.5$ with $\Delta x/N$\\
\code{:dsgs\_Cmin} & $0.06\cdot h_{fac} = 0.075$ & $\equiv C_{\min}=0.06$ with $\Delta x/N$ (\code{JEXPRESSO\_BW\_CMIN})\\
\code{:dsgs\_gamma} & \cm{$=\gamma_{mhd}$} $=2$ & \\
\code{:dsgs\_Prt} & \cm{1.0} & no effect in the conserved form\\
\code{:dsgs\_rel} & \cm{1.0} (default) & floor = the physical scale itself\\
\code{:dsgs\_norms} & \code{"domain"} & means/spreads over the whole tube (MPI-reduced)\\
\code{:dsgs\_conserved} & \mj{true} & $\nu$ on every slot\\
\code{:dsgs\_nazarov\_energy} & \mj{false} (default) & \\
\code{:dsgs\_fast\_floors} & \mj{false} (default) & \\
\code{:dsgs\_thermal\_residual} & \mj{false} (default) & 2D only in any case\\
\code{:dsgs\_hold\_steps} & \mj{2} & \code{JEXPRESSO\_BW\_HOLD}\\
\code{:dsgs\_cutoff} & \cm{0} & \code{JEXPRESSO\_BW\_CUTOFF}\\
\code{:dsgs\_freeze\_stage} & \cm{false} (default) & recomputed at every RK stage\\
\code{:ldsgs\_nodal} & \cm{false} (element form) & \code{JEXPRESSO\_BW\_NODAL}; \code{:dsgs\_Cl => 0.4} only acts in the nodal form\\
\code{:dsgs\_nodal\_rho}, \code{:dsgs\_ref\_weight}, \code{:dsgs\_reference} & \cm{false} & defaults\\
\code{:dsgs\_local\_rel} & \cm{1.0} (default) & unused (domain norms)\\
\end{keytable}

\section{DynSGS as it runs for this case}
\label{sec:dsgs}
Notation: element $K$, LGL node $i$, element lumped mass $m^K_i=\omega_i J_{K,i}$,
\cm{$\Delta_K=\Delta x/(N+1)$} $=1/750$, stage time $t=t^n+\tau$, $h=\Delta t$.

\begin{enumerate}[leftmargin=*]
\item \textbf{History.} \cm{At the first RHS call of every step} ($t-t_{hist}\ge0.999\,\Delta t$) the
\dsgs{} triple is rotated, $(\bq^{n-2},\bq^{n-1},\bq^n)\leftarrow(\bq^{n-1},\bq^n,\bq)$, and the
rotation counter $n_{hist}$ is incremented (\code{rhs.jl:936}). The triple is seeded with the
initial condition.
\item \textbf{Hold.} With $H=2$ the residual right-hand side and the weights are held at zero
while $n_{hist}<\max(2,H)+1=3$, i.e.\ \mj{during the first two steps}; $\nu$ is then the
$C_{\min}$ floor alone (\code{rhs.jl:182, 195, 947}).
\item \textbf{Time derivative} (residual sensor, \code{rhs.jl:964--971}):
\[
D_t q=\begin{cases}
\dfrac{3q^n-4q^{n-1}+q^{n-2}}{2h}, & \tau\le10^{-8}h,\\[8pt]
\dfrac{2\tau+h}{\tau(\tau+h)}\,q(\tau)-\dfrac{\tau+h}{\tau h}\,q^n+\dfrac{\tau}{h(\tau+h)}\,q^{n-1}, & \tau>0 .
\end{cases}
\]
\item \textbf{Residual} of slot $s$ at node $i$ of $K$: $R^s_{K,i}=\big|D_tq^s_i - r^{K,s}_i/m^K_i\big|$,
with $r^K$ the element's own weak inviscid right-hand side (\code{params.rhs\_el}, not assembled).
\textbf{Dirichlet nodes:} the two end nodes are listed by \code{\_dsgs\_boundary\_pairs!} and their
entries replaced by $m^K_i D_tq_i$, so $R\equiv0$ there for all 8 slots (\code{rhs.jl:313, 346}).
\item \textbf{Normalization} (domain): $\langle q^s\rangle$ is the plain average over all element
nodes of the tube (shared nodes counted once per element, MPI-summed), $S^s=\max|q^s-\langle q^s\rangle|$
(MPI-max). With $\bar\rho=\max(|\langle\rho\rangle|,10^{-16})$, $\bar p=(\gamma-1)\max(\langle E\rangle-|\langle\bmom\rangle|^2/(2\bar\rho)-\tfrac12|\langle\bB\rangle|^2,0)$
and the sound speed $\bar c=\sqrt{\gamma\bar p/\bar\rho}$ (no fast floors),
\[
n^s=\max\big(S^s,\ r\,\sigma^s\big)+10^{-16},\qquad
\sigma^\rho=\bar\rho,\ \ \sigma^{\rho u,\rho v,\rho w}=\bar\rho\bar c,\ \ \sigma^E=\bar\rho\bar c^2,\ \ \sigma^{B}=\sqrt{\bar\rho}\,\bar c,
\]
with \cm{$r=$ \code{:dsgs\_rel} $=1$} (in the code: $10^3\cdot$\code{:dsgs\_rel} times the built-in $10^{-3}$).
\item \textbf{Element ratio and cap speed.}
$\mathcal R_K=\max_{i\in K}\max_{s\le 8} R^s_{K,i}/n^s$ (all 8 slots; there is no $\psi$),
$w_K=\max_{i\in K}\big(|\bv_i|+c_{f,i}\big)$ with the \emph{exact fast speed along $x$},
$c_f^2=\tfrac12\big(a^2+b^2+\sqrt{(a^2+b^2)^2-4a^2b_x^2}\big)$, $a^2=\gamma p/\rho$, $b^2=|\bB|^2/\rho$,
$b_x^2=B_x^2/\rho$, $p$ clipped at 0 (\code{SGS.jl:2075--2088}).
\item \textbf{Coefficient} (\code{SGS.jl:2092--2096}):
\[
\nuc=\max\big(0,\min(\Cmax\Delta_Kw_K,\ \CR\Delta_K^2\mathcal R_K)\big),\qquad
\nufl=\Cmin\Delta_Kw_K,\qquad \nu_K=\max(\nuc,\nufl),
\]
\cm{one value per element}, \cm{recomputed at every Runge--Kutta stage}; $(\CR,\Cmax,\Cmin)=(1.5625,0.625,0.075)$
with $\Delta_K=\Delta x/5$, i.e.\ $(1,\,0.5,\,0.06)$ on $\Delta x/4$.
\item \textbf{Slot coefficients} (conserved form): $\mu^s_K=\mu[s]\,\nu_K=\nu_K$ for all 8 slots.
\end{enumerate}

\section{The viscous operator}
\code{user\_primitives.jl} returns the conserved variables themselves, and the 1D loop applies one
scalar flux per slot (\code{rhs.jl:2823}, \code{1349}):
\[
\partial_t\bq+\partial_x\bF(\bq)=\partial_x\big(\nu_K\,\partial_x\bq\big)\quad\text{on }(\rho,\rho u,\rho v,E,\rho w,B_x,B_y,B_z),
\]
\mj{a Laplacian of the conserved state}, in divergence form, with \mj{no separate $\tau\cdot\bv$ or
resistive-work term} (the $E$ Laplacian carries them). $\nu$ is piecewise constant (a staircase
with jumps at element interfaces). $B_x$ is constant in the exact solution; its residual stays in the
max and acts as an oscillation detector.

\section{Positivity}
\begin{itemize}[leftmargin=*]
\item \textbf{No state limiter.} \mj{\code{:lpositivity} is false} (default) and the solver carries
no stage limiter: \mj{nothing modifies the state to enforce $\rho>0$ or $p>0$}.
\item \textbf{Flux-only floor.} \code{pressure\_mhd1d} returns $\max\big((\gamma-1)(E-\dots),\,10^{-9}\big)$
(\code{const p\_floor\_mhd = 1.0e-9}, \code{user\_flux.jl}). It is used by \code{user\_flux!} and by the
output; it changes the flux of a node whose pressure has gone below $10^{-9}$, not the state.
\item \textbf{What keeps the solution admissible} is therefore the dissipation alone: the conserved-form
Laplacian with $\nu\ge\nufl=0.075\,\Delta_Kw_K>0$ everywhere (a background of
$\Cmin/\Cmax=12\%$ of the cap), and the cap $\Cmax\Delta w$ at the discontinuities. Inside \dsgs{} itself the pressure
is clipped at 0 when it builds $\bar p$ and $c_f$, so a negative nodal $p$ cannot produce a NaN coefficient.
\end{itemize}

\section{Observations: code versus comments}
\label{sec:obs}
\begin{enumerate}[leftmargin=*]
\item \textbf{The default (element-form) run should fail on its first \dsgs{} call.} Commit
\code{708bd78} added the keyword \code{lthermal = \dots} to the 1D call
(\code{rhs.jl:1456}) but not to the signature of \code{compute\_dsgs\_viscosity!(::DSGS\_MHD,
::NSD\_1D)} (\code{SGS.jl:1929--1956}, last keyword \code{lfast\_floors}). Julia raises a
\code{MethodError} (``unsupported keyword argument'') for an undeclared keyword. The nodal path
(\code{JEXPRESSO\_BW\_NODAL=1}) does not pass \code{lthermal} and is not affected. This was found by
reading; Julia is not installed in the environment where this note was written, so it was not executed.
The fix is one line: add \code{lthermal::Bool=false} to the 1D signature (it is ignored there).
\item The header of \code{user\_inputs.jl} says ``integrated with SSPRK53''; the deck sets
\code{CarpenterKennedy2N54()}.
\item \code{DSGS.md} \S4.7 says this case needs a 3\% $C_{\min}$ floor; the deck sets
$0.06\,h_{fac}$ (its own comment says 6\%; relative to the cap $\Cmax$ that is 12\%).
\end{enumerate}

% Shared appendix: where each mechanism lives in the source (branch sm/MHD, commit dcc3653).
\section*{Appendix: code map (branch \branch, commit \commit)}
\addcontentsline{toc}{section}{Appendix: code map}
{\small
\begin{longtable}{@{}L{0.47\textwidth}L{0.47\textwidth}@{}}
\toprule
\textbf{Mechanism} & \textbf{Where (file:line)}\\
\midrule\endhead
\bottomrule\endfoot
Deck defaults of every \code{:dsgs\_*} key & \code{src/io/mod\_inputs.jl:1085--1376}\\
Deck defaults of \code{:lpositivity}, floors, report & \code{src/io/mod\_inputs.jl:456--477}\\
\code{:positivity\_method}; rebuild of the RK solver & \code{src/io/mod\_inputs.jl:828--838}\\
\dsgs{} history rotation, stage weights, hold, freeze & \code{src/kernel/operators/rhs.jl:927--984}\\
Stencil selection \code{\_dsgs\_stencil} & \code{rhs.jl:141}\\
Hold length \code{\_dsgs\_hold\_steps} ($\max(2,H)+1$ rotations) & \code{rhs.jl:182}\\
Residual right-hand side (element / legacy / Dirichlet) & \code{rhs.jl:188--280}; legacy fill \code{rhs.jl:282}; boundary zeroing \code{rhs.jl:313, 346}\\
1D \code{DSGS\_MHD} dispatch & \code{rhs.jl:1398--1470}\\
2D \code{DSGS\_MHD} dispatch (sets \code{dsgs\_ref\_weight}, \code{dsgs\_split\_energy}, \code{dsgs\_nodal\_rho}) & \code{rhs.jl:1517--1610}\\
2D viscous assembly \code{\_viscous\_rhs\_el\_2d\_dsgs!} & \code{rhs.jl:1796}\\
2D weak viscous flux \code{\_expansion\_visc!} (conserved flag, ref.\ weight, split energy, $\tau\cdot\bv$, resistive work) & \code{rhs.jl:2943--3210} (\code{lcons} 2959, weight 2978, split 3050, $\tau\cdot\bv$ 3076, $\eta\,\bB\cdot\nabla\bB$ 3091)\\
1D weak viscous flux & \code{rhs.jl:2823}; element loop \code{rhs.jl:1321}\\
2D \code{DSGS\_MHD} element kernel & \code{src/kernel/physics/SGS.jl:1598--1912}\\
1D \code{DSGS\_MHD} element kernel & \code{SGS.jl:1929--2120}\\
Nodal kernels (off in every MHD deck; BW can switch it on with \code{JEXPRESSO\_BW\_NODAL=1}) & \code{SGS.jl:2144} (1D), \code{SGS.jl:2515} (2D)\\
Normalization floor \code{\_dsgs\_denom}, \code{:dsgs\_rel} multiplier & \code{SGS.jl:676--679}; set per call in \code{rhs.jl:192}\\
Smoothness cutoff \code{\_dsgs\_cut} & \code{SGS.jl:703}\\
\code{SGS\_diffusion(::DSGS\_MHD)} (nodal-$\rho$ switch) & \code{SGS.jl:547--569}\\
Element length $\Delta^{elem}_K$ (2D: min of corner extents) & \code{src/kernel/mesh/mesh.jl:6430--6457}; 1D $\Delta x$: \code{mesh.jl:5863}\\
\code{uprimitive} width ($n_{eqs}+2$ in 2D, $n_{eqs}+1$ in 1D) & \code{src/kernel/globalStructs.jl:139--153}\\
Realizability repair (\code{"repair"}, unused by the MHD decks) & \code{src/kernel/positivity/Positivity.jl:294}; driver \code{positivity\_driver.jl:124}; called in \code{rhs.jl:821}\\
Conservative element limiter (\code{"conservative"}) & \code{Positivity.jl:433--559}; driver \code{positivity\_driver.jl:245--337}\\
Main \code{solve} call (\code{saveat} non-empty $\Rightarrow$ \code{calck = true}) & \code{src/kernel/solvers/TimeIntegrators.jl:855}\\
\end{longtable}}
\noindent\small Stage-limiter calls in \code{CarpenterKennedy2N54}: with \code{williamson\_condition = false}
the limiter is called before the right-hand side of stages 2--5 and once after stage 5
(OrdinaryDiffEq v6.90/v6.95, \code{low\_storage\_rk\_perform\_step.jl}). In OrdinaryDiffEq
\code{calck = true} forces \code{williamson\_condition = false} in the 2N cache, and Jexpresso's
\code{solve} passes a non-empty \code{saveat}, so this holds for every deck that hands the solver
a stage limiter. (Jexpresso ships no \code{Manifest.toml}; the package version is whatever the
environment resolves, and the check above was made against v6.90/v6.95.)

\end{document}
